%@ pages 223-232
%@ status translated
% Original section 20, "Exemple: la descente pour les n-champs", pp. 223--232.
%
% Translation notes (the mathematics is left exactly as printed):
%  * "précat (de Segal)" is "precategory", "préchamp (de Segal)" is
%    "prestack", as in sections 2 and 3; "cmf" is rendered "cmc"; "somme
%    disjointe" is "disjoint sum" and "produit direct" is "direct product"
%    (distinct from "produit fibré", "fibre product").
%  * "intérieur $k$-groupique" is "$k$-groupic interior", as in section 2.
%  * In the display after "En effet, on a" the source prints $nSeCHAMP$,
%    without the \underline used everywhere else, and $\Yy /Y$ where
%    $\Xx /Y$ is meant; both are kept as printed (% comments mark the spots).
%  * The source prints "Noua affirmons" for "Nous affirmons"; the prose is
%    translated normally, the mathematics around it is unaffected.

\clearpage
\section{Example: descent for \texorpdfstring{$n$}{n}-stacks}

\label{exemplepage}

The first example of an application of Theorem \ref{dansleschamps} is descent
for $n$-stacks, the analogue of the fact that the prestack of sheaves is a
stack. This descent theorem is classical for $n=1$. For $n=2$ it was stated by
Breen in \cite{Breen}, without proof. Breen uses it to give one of the two
directions of his description of $2$-stacks in terms of cocycles.


\begin{theorem}
\label{descente}
The Segal $n+1$-prestack $nSe\underline{CHAMP}(\Xx )$ is a Segal $n+1$-stack
(i.e.\ $\infty$-stack) over $\Xx$.
\end{theorem}

This theorem means that Segal $n$-stacks glue together, i.e.\ in particular
that descent data for Segal $n$-stacks are effective. However, since we do not
have a very good concrete description of what exactly a descent datum is (see
the end of \S5 for the best that can currently be said), we cannot give a
description, entirely in terms of cocycles, à la Breen \cite{Breen}, of what
this effectivity of descent data means.


\subnumero{Proof using \texorpdfstring{Theorem~\ref{dansleschamps}}{Theorem}}
Suppose that $\Xx$ admits fibre products and has sufficiently many disjoint
sums compatible with fibre products. In that case one can deduce the present
Theorem \ref{descente} directly from Theorem \ref{dansleschamps}.
Indeed, one can verify that the family of cmc's
$$
X\mapsto nSePCh(\Xx /X)
$$
is a left Quillen presheaf, with compatibility with direct products. The
adjoint of the restriction for $s: Y\rightarrow X$ and $A\in
nSePCh(\Xx /Y)$ is defined by
$$
s_{\ast}(A)(U\rightarrow X):= A(Y\times _XU\rightarrow Y).
$$
The restriction evidently preserves cofibrations and trivial cofibrations, and
compatibility with direct products (i.e.\ condition (3) of
\ref{dansleschamps}) follows from an easy computation.

For condition (2),
it suffices to observe that if
$i: U\hookrightarrow X$ is a component
of a disjoint sum, and if $A$ is any kind of presheaf on
$U$, then one has
$i^{\ast}i_{\ast}A=A$.

Condition (1) is an
immediate consequence of the local character of the
notion of $\Gg$-weak equivalence between Segal $n$-prestacks.

The
fibred model categories admit arbitrary limits and colimits, and
functorial factorisations. For the structures of type (II), see
Theorem
\ref{sectionscmf}; as usual we refer to \cite{Jardine} for the
technique needed
to give a rigorous proof of the first part of
the hypothesis of \ref{sectionscmf}
(II). One thus obtains condition (0) of Theorem \ref{dansleschamps}.

%%% verification added in version 2
%%% this is at the top of page 220 in v1

One must verify condition (4)
(added in version 2 of the paper) of Theorem \ref{dansleschamps}.
So let $f:X\rightarrow
Y$ be a morphism in $\Xx$. The point is to see that
$$
f^{\ast} :L(nSePCh , W)(Y) \rightarrow L(nSePCh , W) (X)
$$
preserves (homotopy) limits. Here $f^{\ast}$ is just the restriction
to $\Xx /X$ of the Segal $n$-stacks on $\Xx /Y$. Limits of
Segal $n$-stacks are computed object by object, so the restriction preserves
these limits. The fact that limits are computed object by object
is a generalisation of Lemma \ref{fiprodchamp}, which concerns the case of the fibre
product. For the general case one can use the commutation of limits
(\cite{Vogt2}, \cite{limits} 3.4.11) as well as criterion \ref{critereaveclim} (c) to see that a
limit (as a prestack)
of stacks is still a stack. Limits of prestacks are computed
object by object according to \cite{limits} 3.4.4.

%%% end of the v2 verification; I also changed the first sentence of the following.

%%%%% for v3 we verify condition (5) by modifying slightly the
%%%%%% wording of the following paragraph %%%%%

It will then only remain to verify hypothesis (5) of
\ref{dansleschamps} (new in {\tt v3})
in order to prove that $L(nSePCh, W)$
is a Segal $1$-stack. We shall do this at the same time as the proof, by
criterion \ref{critere}, that $nSe\underline{CHAMP}$ is a stack.

Recall that by \ref{intereqloc}, $L(nSePCh, W)$
is equivalent (for the coarse topology)
to the $1$-groupic interior
$$
L(nSePCh, W) \cong nSe\underline{CHAMP}^{int,1}.
$$
Part (a) of criterion \ref{critere} for $nSe\underline{CHAMP}$
is automatic because, by definition, the
objects of $nSe\underline{CHAMP}(X)$ are the $\Gg$-fibrant Segal $n$-stacks,
and if $A$ and $B$ are $\Gg$-fibrant on $\Xx /X$
then
$$
nSe\underline{CHAMP}_{1/}(A,B) :=
\underline{Hom}(A,B)
$$
is fibrant---in particular it is a stack.
It follows that the $0$-groupic interior of the latter,
$nSe\underline{CHAMP}_{1/}(A,B)$, is a stack. The object of morphisms in
the $1$-groupic interior is also the $0$-groupic interior of
the object of morphisms, so we obtain that
the $1$-groupic interior
$nSe\underline{CHAMP}^{int,1}$
(equivalent to $L(nSePCh, W)$)
satisfies part (a) of criterion \ref{critere}. Now this condition
is exactly condition (5) of Theorem \ref{dansleschamps}.

We have therefore finished verifying the hypotheses of \ref{dansleschamps},
and by this theorem the localised category $L(nSePCh, W)$
is a Segal $1$-stack.
From this result one obtains part (b) of the criterion of \ref{critere}
(see also \ref{devissage}) for
$nSe\underline{CHAMP}$.

Proposition \ref{critere}
now applies to conclude that
$nSe\underline{CHAMP}$ is
a Segal $n+1$-stack.
\eop

%%%%%%%

\subnumero{A direct proof}
We now give a direct proof of Theorem \ref{descente},
which does not use
Theorem \ref{dansleschamps}. This proof seems
to give more indications than the first one about the way in which
Segal $n$-stacks glue together. Moreover, it applies
without any further hypothesis on
$\Xx$ and adapts {\em mutatis mutandis} to the case of $n$-stacks that are not
of Segal.

In what follows we fix
a site $\Xx$, which we shall no longer mention
in the notation. When needed we shall write
$\Gg$ for the
topology on $\Xx$ (as opposed to the coarse topology).

We shall use the construction $\Upsilon$ (or its variants $\Upsilon ^k$) of
\cite{limits}. It generalises immediately
to Segal $n$-precategories and, by functoriality, to presheaves of
Segal $n$-precategories, i.e.\ to Segal $n$-prestacks, while preserving its
universal properties, of which we shall sometimes have need.


Let
$$
nSe\underline{CHAMP}\rightarrow nSe\underline{CHAMP}'
\rightarrow nSe\underline{CHAMP}''
$$
denote two cofibrations: the first is a trivial cofibration for the coarse
topology (in particular an object-by-object equivalence) towards a
fibrant model for the coarse topology; the second is a trivial cofibration
for the topology $\Gg$ towards a fibrant model for the topology $\Gg$. Their
composite is likewise a trivial cofibration for the topology
$\Gg$ towards
a fibrant model for the topology $\Gg$.

By construction the $nSe\underline{CHAMP}(X)$ are already
Segal $n+1$-categories. The condition that $nSe\underline{CHAMP}$ be
a Segal $n+1$-stack is therefore equivalent to the condition that the morphism
$$
nSe\underline{CHAMP}'\rightarrow nSe\underline{CHAMP}''
$$
be an equivalence for the coarse topology, i.e.\ an
object-by-object equivalence.

We shall use criterion \ref{critere}.
For $A$ and $B$ in $nSe\underline{CHAMP}(X)$
one has by construction
$$
nSe\underline{CHAMP}(X)_{1/} (A,B) = \underline{Hom}(A,B).
$$
In particular, since $A$ and $B$ are cofibrant and fibrant (by
the very definition of $nSe\underline{CHAMP}$), this is a Segal $n$-stack.
The first hypothesis of \ref{critere} is therefore verified.


It follows that, for every $X$, the morphism
$$
nSe\underline{CHAMP}(X)
\rightarrow
nSe\underline{CHAMP}''(X)
$$
is fully faithful. Indeed, the morphism
$$
nSe\underline{CHAMP}
\rightarrow
nSe\underline{CHAMP}''
$$
is a $\Gg$-weak equivalence of Segal $n+1$-prestacks whose values are Segal
$n+1$-categories; by definition, for every $X$ and all
$A,B\in nSe\underline{CHAMP}'(X)_0$, the morphism
$$
nSe\underline{CHAMP}_{1/}(A,B)
\rightarrow
nSe\underline{CHAMP}''_{1/}(A,B)
$$
is a weak equivalence of Segal $n$-prestacks whose
source and target
are already Segal $n$-stacks; by \ref{morphentrechamps}
it is a weak equivalence for the coarse topology.
This result admits the following generalisation:

\begin{lemma}
\label{hompleinfidele}
If $A$ is a Segal $n+1$-prestack then the morphism of
Segal $n+1$-prestacks on $\Xx$
$$
\underline{Hom}(A, nSe\underline{CHAMP}')\rightarrow
\underline{Hom}(A,nSe\underline{CHAMP}'')
$$
is fully faithful for the coarse topology.
\end{lemma}
{\em Proof:}
Let $Z$ be the full subprestack of $nSe\underline{CHAMP}''$
consisting of the objects that are in the image of $nSe\underline{CHAMP}'$.
Then the morphism
$nSe\underline{CHAMP}'\rightarrow Z$ is a weak equivalence for the
coarse topology. Hence the same is true of the morphism
$$
\underline{Hom}(A, nSe\underline{CHAMP}')\rightarrow
\underline{Hom}(A, Z).
$$
But the morphism
$$
\underline{Hom}(A, Z)\rightarrow
\underline{Hom}(A, nSe\underline{CHAMP}'')
$$
is fully faithful for the coarse topology: if $a,b$ are
two objects of $\underline{Hom}(A, Z)(X)$ then a morphism
$$
U\rightarrow \underline{Hom}(A, Z)_{1/}(a,b)
$$
corresponds to a morphism
$$
A\times \Upsilon (U)\rightarrow Z
$$
which gives $a$ (resp.\ $b$) on $A\times 0$ (resp.\ $A\times 1$);
since all the objects of $A\times \Upsilon (U)$ are either in
$A\times 0$ or in $A\times 1$, this
still corresponds to a morphism
$$
A\times \Upsilon (U)\rightarrow nSe\underline{CHAMP}''
$$
which gives $a$ (resp.\ $b$) on $A\times 0$ (resp.\ $A\times 1$),
in other words to a morphism
$$
U\rightarrow \underline{Hom}(A, nSe\underline{CHAMP}'')_{1/}(a,b).
$$
In fact, one has the equality
$$
\underline{Hom}(A, Z)_{1/}(a,b)  =
\underline{Hom}(A, nSe\underline{CHAMP}'')_{1/}(a,b).
$$
By composition, one obtains that
$$
\underline{Hom}(A, nSe\underline{CHAMP}')\rightarrow
\underline{Hom}(A, nSe\underline{CHAMP}'')
$$
is fully faithful for the coarse topology.
\eop


We continue the proof of Theorem \ref{descente}.
In view of the criterion of Proposition \ref{critere},
it is now a matter of proving that for every covering
sieve $\Bb \subset \Xx /X$, every
morphism
$$
f:\ast _{\Bb} \rightarrow nSe\underline{CHAMP}'|_{\Xx /X}
$$
extends to a morphism defined on $\ast _X$. Since the formation of
$nSe\underline{CHAMP}$ is compatible with base change, and since
the fibrant replacement for the coarse topology can also be
chosen compatible with restrictions (Lemma \ref{shriek}), we may place
ourselves henceforth on the site $\Xx /X$ and delete from our notation the
reference to the ``restriction
to $\Xx
/X$''.

Let $$
E:= \underline{Hom}(\ast _{\Bb},nSe\underline{CHAMP}')_{1/}(\ast , f).
$$
Here $\underline{Hom}$ is the internal $Hom$ of Segal $n+1$-prestacks on
$\Xx /X$. In particular, $E$ is a Segal $n$-prestack on $\Xx /X$.

The morphism
$$
\underline{Hom}(\ast _{\Bb},nSe\underline{CHAMP}')_{1/}(\ast , f)
\rightarrow
\underline{Hom}(\ast _{\Bb},nSe\underline{CHAMP}'')_{1/}(\ast , f)
$$
is a weak equivalence of $n$-prestacks
for the coarse topology on
$\Xx /X$ by Lemma \ref{hompleinfidele}. Since the target of this
morphism is
$\Gg$-fibrant, it is a Segal $n$-stack, hence
its source
---which is $E$---is
also a Segal $n$-stack. The latter being also fibrant for the
coarse topology, we obtain that $E$ is $\Gg$-fibrant by
\ref{grofibchampimplfib}.

The idea of the proof is that $E$ solves the descent problem that is
posed. A little more technique is needed
to obtain a morphism
between $f$ and the constant morphism $cst(E)$, and
to complete the proof of the theorem.

We first define a Segal $n+1$-prestack $\Upsilon (E)$ by
applying the
definition of \cite{limits} object by object, i.e.
$$
\Upsilon (E)(Y):= \Upsilon (E(Y))
$$
(the extension of the constructions of \cite{limits} object by object will be used
below without further remark).

By the universal
property that defines $\Upsilon$ one obtains a morphism $$
\Upsilon (E)\rightarrow \underline{Hom}(\ast _{\Bb},nSe\underline{CHAMP}')
$$
or equivalently
$$
h: \ast _{\Bb} \times \Upsilon (E) \rightarrow nSe\underline{CHAMP}'.
$$
One has
$$
h|_{\ast _{\Bb} \times 0} = \ast
$$
and
$$
h|_{\ast _{\Bb} \times 1} = f.
$$

We now use the extension properties for morphisms defined
on parts of $\Upsilon ^k$---we
refer to \cite{limits} for the details---,
noting that they extend to the framework of prestacks.
Recall the notation $01$, $12$ and $02$ for the copies of $\Upsilon (A)$,
$\Upsilon (B)$ and $\Upsilon (A\times B)$ respectively in $\Upsilon ^2(A,B)$.

By the analogue of \cite{limits} 4.3.5, there exists a morphism
$$
h^2: \ast _{\Bb} \times \Upsilon ^2(E,\ast )
$$
with
$$
h^2|_{01} = 1_E,
$$
and
$$
h^2|_{02} = h.
$$
Recall that in the notation of \cite{limits}
$1_E$ is a map
$$
\Upsilon (E) \rightarrow nSe\underline{CHAMP}'
$$
which induces $\ast$ on the first vertex $0$; and
which induces the constant morphism
$$
cst(E) :\ast \rightarrow nSe\underline{CHAMP}'
$$
with value $E$, on the second vertex $1$.

The restriction of $h$ to the edge $12$ therefore provides a morphism
$$
g: \ast _{\Bb} \times I \rightarrow nSe\underline{CHAMP}',
$$
where we have written $I$ for the $1$-category with two objects and one arrow,
which is also equal to $\Upsilon (\ast )$.
The objects of $I$ will be called $0$ and $1$ although this notation is
incompatible with the notation for the vertices of $\Upsilon ^2(E,\ast )$
(the object $0$ of $I$ corresponding to the vertex $1$,
and the object $1$ of $I$ corresponding to the
vertex $2$).

The image of the first object $0$ of $I$ (hence the image of the vertex $1$ of
$\Upsilon ^2$)
is the constant map $cst(E)$ composed with the projection $\ast
_{\Bb}\rightarrow \ast$. The image
of the second object $1$ of $I$ is the map $f$.

We assert that $g$ extends to a morphism
$$
\overline{g}: \ast _{\Bb} \times \overline{I} \rightarrow
nSe\underline{CHAMP}'
$$
where $\overline{I}$ is the $1$-category with two objects $0,1$ and an
isomorphism between them (one has $I\subset \overline{I}$). One
first observes the following property:

$(\ast
)$ for every
$Y\in \Bb$, the morphism induced by $g$
$$
E(Y)\rightarrow f_Y(\ast _{\Bb}(Y))
$$
is a weak equivalence of Segal $n+1$-precategories.

Indeed, one has
$$
E(Y)= \Gamma \underline{Hom}(\ast _{\Bb }|_{\Xx /Y},
nSe\underline{CHAMP}'|_{\Xx /Y})_{1/}(\ast , f),
$$
but also $\ast _{\Bb }|_{\Xx /Y}=\ast _{\Xx /Y}$ and therefore
$$
E(Y)= \Gamma (\Yy /Y, nSe\underline{CHAMP}')_{1/}(\ast , f|_{\Xx /Y})
% the source prints \Yy/Y here (for \Xx/Y) and nSeCHAMP without \underline
$$
$$
\cong nSeCHAMP (Y)_{1/}(\ast , f_Y(\ast _{\Bb}(Y))) \cong
f_Y(\ast _{\Bb}(Y)).
$$

Let $C$ denote the result obtained
from $\ast _{\Bb} \times I$, by freely adding
above each object of $\Bb$ a cell corresponding to the cofibration
$I\hookrightarrow \overline{I}$. One has a morphism from $C$ to
$\ast_{\Bb} \times \overline{I}$. This morphism is an equivalence, because by
\cite{limits} 2.5.1, doing (in order to invert the same arrow)
the coproduct with the cofibration $I\rightarrow
\overline{I}$ several times has the same effect as doing it only once.
Now the statement $(\ast )$ above implies that $g$ extends to a
morphism $C\rightarrow nSe\underline{CHAMP}'$. By the usual arguments,
this implies that $g$ extends to $\overline{g}$ as announced.

Thanks to this result, one can set
$$
V:= (\ast _{\Bb}\times \overline{I}) \cup ^{\ast _{\Bb} \times 0}\ast _X.
$$
The map $V\rightarrow \ast _X$ is a weak equivalence for the
coarse topology. The map
$$
cst(E): \ast _X \rightarrow nSe\underline{CHAMP}'
$$
glues with the map
$$
\overline{g}: \ast _{\Bb} \times \overline{I} \rightarrow
nSe\underline{CHAMP}'
$$
because one has, by construction,
$$
\overline{g}|_{\ast _{\Bb} \times 0} = cst(E).
$$
One obtains a map
$$
V\rightarrow nSe\underline{CHAMP}'
$$
whose restriction to $\ast _{\Bb} \times 1$ is equal to $f$.
Let
$$
V\rightarrow V' \rightarrow \ast _X
$$
be a factorisation with a trivial cofibration as first morphism and
a trivial fibration as
second morphism (all of this for the coarse topology).
The above map extends to a map
$$
\alpha : V'\rightarrow nSe\underline{CHAMP}'
$$
whose restriction to
$\ast _{\Bb} \times 1\subset V \subset V'$ is still equal to $f$.
Consider the diagram
$$
\begin{array}{ccc}
\ast _{\Bb} & \rightarrow & V' \\
\downarrow && \downarrow \\
\ast _X & = & \ast _X.
\end{array}
$$
The top arrow is the inclusion
$$
i: \ast _{\Bb}=\ast _{\Bb} \times 1\subset V
\subset V'.
$$
One can write $\alpha \circ i = f$.
The right-hand arrow is a trivial fibration for the coarse topology,
in particular it has the lifting property for every
cofibration.
Hence there exists a lift $j:\ast _X\rightarrow V'$ equal to
$i$ on $\ast _{\Bb}$.   The morphism
$$
\alpha \circ j: \ast _X \rightarrow nSe\underline{CHAMP}'
$$
is the extension of $f$ that we are looking for, which completes the
proof
of Theorem \ref{descente}.
\eop



\begin{corollary}
The presheaves of simplicial categories $X\mapsto L(nSePCh(\Xx /X))$
are Segal $1$-stacks. In particular the presheaf of simplicial
categories
$$
X\mapsto L(PrefSpl(\Xx /X),W^{\rm Illusie})
$$
is a Segal $1$-stack.
\end{corollary}
{\em Proof:}
By Theorem \ref{descente}, $nSe\underline{CHAMP}$ is a Segal $n+1$-stack.
Hence (cf.\ \S10) its $1$-groupic interior is a Segal $1$-stack.
One applies \ref{intereqloc}.
\eop

This corollary shows that the homotopy descent data for simplicial
presheaves or Segal $n$-prestacks are effective, up to
equivalence.


Let $nSeChamp(\Xx )$ be the strict $1$-category of presheaves of
Segal $n$-categories that are Segal $n$-stacks on $\Xx$.
Here the notion of $\Gg$-equivalence coincides with that of object-by-object
equivalence. One obtains that {\em for
$L(nSeChamp(\Xx ))$ the descent data are effective}, which means
that Segal $n$-stacks glue together, up to object-by-object
equivalence. This statement is
the analogue of the classical gluing result for sheaves.

\subnumero{Statements for \texorpdfstring{$n$}{n}-stacks (not of Segal)}

Essentially, all the statements of the
present work, except those concerning
directly the Dwyer--Kan localised categories, are also valid for
$n$-categories, $n$-stacks etc., not of Segal. Without entering into details,
we indicate here the main facts (among other things, in order to justify the
title chosen!).

First, the category $nPC$ of $n$-precategories replaces $nSePC$, that
of Segal $n$-precategories, and one has at one's disposal the operation $Cat$ (analogous to
$SeCat$), and the notion of $n$-category has the same definition---all of this
is to be found in
\cite{Tamsamani} and \cite{nCAT}.  An {\em $n$-prestack on $\Xx$} is a
presheaf on $\Xx$ with values in $nPC$; we denote by
$nPCh(\Xx )$ the category of
$n$-prestacks.

The categories $nPC$ and $nPCh(\Xx )$ have
cmc structures analogous to those of \ref{SeCmf} and \ref{cmf}, and
$nPCh(\Xx )$ also has the structure of type HBKQ of \ref{cmfHirschho}.

One says (in the same way as in \S9) that an $n$-prestack $A$ is
an {\em
$n$-stack} if the morphism $A\rightarrow A'$ of fibrant replacement for the
topology $\Gg$ is an object-by-object equivalence. The properties of
\S9 and the other definition of \S10 adapt {\em mutatis mutandis}.

The cmc's $nPC$ and $nPCh(\Xx )$ are internal (\S11; for $nPC$ this is in
\cite{nCAT}), which allows one to define the $n+1$-categories
$$
nCAT\;\;\; \mbox{and} \;\;\; nCHAMP(\Xx )
$$
as well as the $n+1$-stack $n\underline{CHAMP}(\Xx )$.

One defines a universal family and a morphism $\Phi$ as in \S12.
Theorem \ref{correlation} takes the following form.


\begin{theorem}
\label{correlationNdS}
Let $nCAT'$ be a fibrant replacement of $nCAT$.
If the topology of $\Xx$ is
coarse then $nCHAMP(\Xx )$ is equivalent via $\Phi$ to the
$n+1$-category $\underline{Hom}(\Xx ^o, nCAT')$.

If $\Xx$ is equipped with an arbitrary topology
$\Gg$ then $nCHAMP(\Xx )$ is equivalent to the full subcategory
of
$\underline{Hom}(\Xx ^o, nCAT')$ consisting of the morphisms
$F:\Xx ^o\rightarrow nCAT'$ that satisfy the descent condition
(\cite{limits} 6.3), which says that the arrow
$$
\lim _{\leftarrow} F|_{\Xx /X} \rightarrow
\lim _{\leftarrow} F|_{\Bb}
$$
is an equivalence for every sieve $\Bb \subset \Xx /X$ of $\Gg$.
\end{theorem}
\eop

One has an ``associated stack'' functor between $n+1$-categories
$$
{\bf ch} : nCHAMP (\Xx ^{\rm gro})\rightarrow nCHAMP (\Xx ^{\Gg})
$$
with natural transformation $F\rightarrow {\bf ch}(F)$,
and (Theorem \ref{cestadjoint}) this natural
transformation makes this functor the homotopy adjoint of the inclusion
$$
nCHAMP (\Xx ^{\Gg})\subset nCHAMP (\Xx ^{\rm gro}).
$$

Finally, the main theorem of the present chapter \ref{descente} becomes:

\begin{theorem}
\label{descenteNdS}
The $n+1$-prestack  $n\underline{CHAMP}(\Xx )$ is an $n+1$-stack
over $\Xx$.
\end{theorem}

The direct proof adapts {\em mutatis mutandis}. If one wants to obtain a proof
by applying \ref{dansleschamps}, one must observe the following fact:

\begin{proposition}
\label{segalounon}
The $n+1$-category $nCAT$, considered
by induction as a Segal $n+1$-category, is equivalent
to the full Segal sub-$n+1$-category of $nSeCAT$ consisting
of the objects that are $n$-truncated.

The $n+1$-category $nCHAMP(\Xx )$,
considered in the same way as a Segal $n+1$-category,
is equivalent to the full Segal sub-$n+1$-category
of
$nSeCHAMP$ consisting of the Segal $n$-stacks that are $n$-truncated
object by object over $\Xx$.

The Segal $n+1$-prestack $n\underline{CHAMP}(\Xx )$,
considered by induction as a Segal $n+1$-prestack,
is equivalent to the full Segal subprestack of
$nSe\underline{CHAMP}(\Xx )$
consisting of the objects whose values (on every $X\in \Xx$) are
Segal $n$-stacks (on $\Xx/X$) that are $n$-truncated object by object.
\end{proposition}
\eop

One can then apply \ref{preservetronque}.
\eop
